\documentclass[12pt,a4paper]{article}

%================= PAQUETES =================
\usepackage[utf8]{inputenc}
\usepackage[T1]{fontenc}
\usepackage[spanish,es-noshorthands,es-tabla]{babel}
\usepackage{geometry}
\usepackage{amsmath}
\usepackage{graphicx}
\usepackage{xcolor}
\usepackage{booktabs}
\usepackage{enumitem}
\usepackage{listings}
\usepackage[most]{tcolorbox}
\usepackage{hyperref}
\usepackage{tabularx}

%================= CONFIGURACIÓN =================
\geometry{a4paper, margin=2.5cm}

\definecolor{azulTitulo}{RGB}{30,60,120}
\definecolor{verdePython}{RGB}{53,114,53}
\definecolor{codigoFondo}{RGB}{245,247,250}
\definecolor{codigoComentario}{RGB}{140,140,140}
\definecolor{codigoCadena}{RGB}{170,60,30}
\definecolor{codigoKeyword}{RGB}{0,80,160}
\definecolor{codigoNumero}{RGB}{150,30,150}

\hypersetup{
    colorlinks=true,
    linkcolor=azulTitulo,
    urlcolor=cyan,
    pdftitle={Solucionario - Curso de Programación en Python para Secundaria},
}

% Configuración de listings para bloques de código Python
\lstdefinestyle{pythonStyle}{
    language=Python,
    backgroundcolor=\color{codigoFondo},
    basicstyle=\ttfamily\small,
    keywordstyle=\color{codigoKeyword}\bfseries,
    commentstyle=\color{codigoComentario}\itshape,
    stringstyle=\color{codigoCadena},
    numberstyle=\tiny\color{codigoComentario},
    numbers=left,
    numbersep=8pt,
    showstringspaces=false,
    breaklines=true,
    breakatwhitespace=false,
    tabsize=4,
    frame=single,
    rulecolor=\color{black!20},
    captionpos=b,
    literate=
      {á}{{\'a}}1 {é}{{\'e}}1 {í}{{\'i}}1 {ó}{{\'o}}1 {ú}{{\'u}}1
      {Á}{{\'A}}1 {É}{{\'E}}1 {Í}{{\'I}}1 {Ó}{{\'O}}1 {Ú}{{\'U}}1
      {ñ}{{\~n}}1 {Ñ}{{\~N}}1 {ü}{{\"u}}1 {Ü}{{\"U}}1
      {¿}{{?`}}1 {¡}{{!`}}1
}

\newtcolorbox{cajaSolucion}{
    colback=green!5!white,
    colframe=verdePython,
    boxrule=0.8pt,
    arc=2mm,
    fonttitle=\small,
    left=6pt, right=6pt, top=4pt, bottom=4pt
}

\newtcolorbox{cajaNota}{
    colback=orange!5!white,
    colframe=orange!70!black,
    boxrule=0.8pt,
    arc=2mm,
    title={\textbf{Nota pedagógica}},
    fonttitle=\small,
    left=6pt, right=6pt, top=4pt, bottom=4pt
}

\title{\textbf{Solucionario}\\[6pt]
\large Curso de Programación en Python para Secundaria}
\author{Prof. César Rodríguez}
\date{\today}

\begin{document}

\maketitle

\begin{abstract}
\noindent Este documento contiene las soluciones de todos los ejercicios y retos propuestos en el curso. Está dirigido al docente. Se recomienda que el estudiante intente resolver cada ejercicio \textbf{antes} de revisar la solución. Las soluciones están organizadas por módulo y numeradas igual que en el documento principal.
\end{abstract}

\tableofcontents
\newpage

% ============================================================
% MÓDULO 1
% ============================================================
\section{Soluciones del Módulo 1}

\begin{cajaSolucion}
\textbf{Ejercicio 1}
\begin{lstlisting}[style=pythonStyle]
print("Mi nombre es Ana López")
\end{lstlisting}
\end{cajaSolucion}

\begin{cajaSolucion}
\textbf{Ejercicio 2}
\begin{lstlisting}[style=pythonStyle]
print("Mi nombre es Ana López")
print("Soy de La Paz, Bolivia")
print("Mi hobby es la programación")
\end{lstlisting}
\end{cajaSolucion}

\begin{cajaSolucion}
\textbf{Ejercicio 3}
\begin{lstlisting}[style=pythonStyle]
print("*****")
print("*****")
print("*****")
\end{lstlisting}
\end{cajaSolucion}

\begin{cajaSolucion}
\textbf{Ejercicio 4}

Una posible respuesta:
\begin{lstlisting}[style=pythonStyle]
# Programa: hola mundo
# Autor: Ana López
# Fecha: 2026

# La funcion print muestra un mensaje en pantalla
print("Hola, mundo!")  # Este es el mensaje que se mostrara
\end{lstlisting}
\end{cajaSolucion}

\begin{cajaSolucion}
\textbf{Reto 1}

Cuando \texttt{print()} recibe dos argumentos separados por coma, los muestra en la misma línea separados por un espacio:
\begin{lstlisting}[style=pythonStyle]
# print con varios argumentos los separa automaticamente con un espacio
print("Hola", "mundo")  # salida: Hola mundo

# Programa con mi nombre y mi edad
print("Mi nombre es", "Ana López")
print("Tengo", 15, "años")
\end{lstlisting}
\textbf{Salida:}
\begin{lstlisting}[style=pythonStyle]
Hola mundo
Mi nombre es Ana López
Tengo 15 años
\end{lstlisting}
\end{cajaSolucion}

% ============================================================
% MÓDULO 2
% ============================================================
\section{Soluciones del Módulo 2}

\begin{cajaSolucion}
\textbf{Ejercicio 5}
\begin{lstlisting}[style=pythonStyle]
edad = 15
nombre = "Ana"
estatura = 1.62

print(nombre)
print(edad)
print(estatura)

print(type(nombre))       # <class 'str'>
print(type(edad))         # <class 'int'>
print(type(estatura))     # <class 'float'>
\end{lstlisting}
\end{cajaSolucion}

\begin{cajaSolucion}
\textbf{Ejercicio 6}
\begin{lstlisting}[style=pythonStyle]
base = 8
altura = 5
area = base * altura

print("El area del rectangulo es:", area)   # 40
\end{lstlisting}
\end{cajaSolucion}

\begin{cajaSolucion}
\textbf{Ejercicio 7}

El error es que se intenta sumar texto (``\texttt{y tengo}'') con un entero (\texttt{edad}), y Python no permite sumar \texttt{str} con \texttt{int}. Hay que convertir \texttt{edad} a texto con \texttt{str()}:
\begin{lstlisting}[style=pythonStyle]
nombre = "Carlos"
edad = 14
# Corregido: convertimos edad a texto con str()
print("Hola, me llamo " + nombre + " y tengo " + str(edad) + " anios")
\end{lstlisting}

\begin{cajaNota}
También se puede resolver con \texttt{print("Hola, me llamo", nombre, "y tengo", edad, "anios")}, donde \texttt{print()} convierte automáticamente cada valor a texto. Esta es la alternativa más didáctica.
\end{cajaNota}
\end{cajaSolucion}

\begin{cajaSolucion}
\textbf{Ejercicio 8}
\begin{lstlisting}[style=pythonStyle]
a = 5
b = 3

print("Antes de intercambiar: a =", a, "b =", b)

# Usamos una variable auxiliar (una tercera caja)
temporal = a
a = b
b = temporal

print("Despues de intercambiar: a =", a, "b =", b)
\end{lstlisting}
\end{cajaSolucion}

\begin{cajaSolucion}
\textbf{Reto 2}
\begin{lstlisting}[style=pythonStyle]
radio = 7
pi = 3.14159
perimetro = 2 * pi * radio

print("El perimetro es", perimetro)    # El perimetro es 43.98226
\end{lstlisting}
\end{cajaSolucion}

% ============================================================
% MÓDULO 3
% ============================================================
\section{Soluciones del Módulo 3}

\begin{cajaSolucion}
\textbf{Ejercicio 9}
\begin{lstlisting}[style=pythonStyle]
nombre = input("Escribe tu nombre: ")
edad = int(input("Escribe tu edad: "))

print("Hola,", nombre + ". El proximo anio tendras", edad + 1, "anios.")
\end{lstlisting}
\end{cajaSolucion}

\begin{cajaSolucion}
\textbf{Ejercicio 10}
\begin{lstlisting}[style=pythonStyle]
a = int(input("Escribe el primer numero: "))
b = int(input("Escribe el segundo numero: "))

print(a, "+", b, "=", a + b)
print(a, "-", b, "=", a - b)
print(a, "*", b, "=", a * b)
print(a, "/", b, "=", a / b)
\end{lstlisting}
\end{cajaSolucion}

\begin{cajaSolucion}
\textbf{Ejercicio 11}
\begin{lstlisting}[style=pythonStyle]
base = float(input("Escribe la base del triangulo: "))
altura = float(input("Escribe la altura del triangulo: "))

area = (base * altura) / 2

print("El area del triangulo es:", area)
\end{lstlisting}
\end{cajaSolucion}

\begin{cajaSolucion}
\textbf{Ejercicio 12}
\begin{lstlisting}[style=pythonStyle]
precio = float(input("Escribe el precio del producto: "))

iva = precio * 0.12
precio_final = precio + iva

print("Precio original:", precio)
print("Monto del IVA (12%):", iva)
print("Precio final:", precio_final)
\end{lstlisting}
\end{cajaSolucion}

\begin{cajaSolucion}
\textbf{Reto 3}
\begin{lstlisting}[style=pythonStyle]
segundos = int(input("Escribe una cantidad de segundos: "))

horas = segundos // 3600
resto = segundos % 3600
minutos = resto // 60
seg = resto % 60

print(segundos, "segundos son", horas, "hora(s),",
      minutos, "minuto(s) y", seg, "segundo(s)")
\end{lstlisting}
\textbf{Prueba con 10000:} 10000 segundos son 2 horas, 46 minutos y 40 segundos.
\end{cajaSolucion}

% ============================================================
% MÓDULO 4
% ============================================================
\section{Soluciones del Módulo 4}

\begin{cajaSolucion}
\textbf{Ejercicio 13}

Resultados esperados:
\begin{lstlisting}[style=pythonStyle]
print(15 // 4)      # 3
print(15 % 4)       # 3
print(2 ** 10)      # 1024
print(10 - 2 * 3)   # 4
print((10 - 2) * 3) # 24
print(9 % 2)        # 1
\end{lstlisting}

\begin{cajaNota}
Este ejercicio practica la \textbf{precedencia de operadores}: la multiplicación se evalúa antes que la suma, pero los paréntesis tienen prioridad total.
\end{cajaNota}
\end{cajaSolucion}

\begin{cajaSolucion}
\textbf{Ejercicio 14}
\begin{lstlisting}[style=pythonStyle]
numero = int(input("Escribe un numero: "))

if numero % 2 == 0:
    print("El numero es par")
else:
    print("El numero es impar")
\end{lstlisting}
\end{cajaSolucion}

\begin{cajaSolucion}
\textbf{Ejercicio 15}
\begin{lstlisting}[style=pythonStyle]
nota1 = float(input("Escribe la primera nota: "))
nota2 = float(input("Escribe la segunda nota: "))
nota3 = float(input("Escribe la tercera nota: "))

promedio = (nota1 + nota2 + nota3) / 3

print("El promedio es:", round(promedio, 2))
\end{lstlisting}
\end{cajaSolucion}

\begin{cajaSolucion}
\textbf{Ejercicio 16}
\begin{lstlisting}[style=pythonStyle]
numero = int(input("Escribe un numero de 4 cifras: "))

miles = numero // 1000
centenas = (numero % 1000) // 100
decenas = (numero % 100) // 10
unidades = numero % 10

print("Digitos:", miles, centenas, decenas, unidades)
\end{lstlisting}
\end{cajaSolucion}

\begin{cajaSolucion}
\textbf{Reto 4}
\begin{lstlisting}[style=pythonStyle]
numero = int(input("Escribe un numero de 3 cifras: "))

centenas = numero // 100
unidades = numero % 10

if centenas == unidades:
    print(numero, "es un numero capicua")
else:
    print(numero, "no es un numero capicua")
\end{lstlisting}

\begin{cajaNota}
Basta comparar centenas con unidades. Si coinciden, el número se lee igual al derecho y al revés (la cifra del medio no afecta). Ejemplos: 121 sí es capicúa, 123 no lo es.
\end{cajaNota}
\end{cajaSolucion}

% ============================================================
% MÓDULO 5
% ============================================================
\section{Soluciones del Módulo 5}

\begin{cajaSolucion}
\textbf{Ejercicio 17}
\begin{lstlisting}[style=pythonStyle]
numero = int(input("Escribe un numero: "))

if numero > 0:
    print("El numero es positivo")
elif numero < 0:
    print("El numero es negativo")
else:
    print("El numero es cero")
\end{lstlisting}
\end{cajaSolucion}

\begin{cajaSolucion}
\textbf{Ejercicio 18}
\begin{lstlisting}[style=pythonStyle]
numero = int(input("Escribe un numero entero: "))

if numero % 2 == 0:
    print("El numero es par")
else:
    print("El numero es impar")

# Condiciones anidadas para el segundo dato
if numero > 100:
    print("El numero es mayor que 100")
else:
    print("El numero es menor o igual que 100")
\end{lstlisting}

\begin{cajaNota}
Las dos comprobaciones son \textbf{independientes}, por eso se resuelve cómodo con dos bloques \texttt{if} secuenciales. También es válido anidarlas dentro de un solo \texttt{if}, como pide el enunciado; ambas formas son correctas.
\end{cajaNota}
\end{cajaSolucion}

\begin{cajaSolucion}
\textbf{Ejercicio 19}
\begin{lstlisting}[style=pythonStyle]
a = int(input("Escribe el primer numero: "))
b = int(input("Escribe el segundo numero: "))

if a > b:
    print("El mayor es:", a)
elif b > a:
    print("El mayor es:", b)
else:
    print("Los numeros son iguales")
\end{lstlisting}
\end{cajaSolucion}

\begin{cajaSolucion}
\textbf{Ejercicio 20}
\begin{lstlisting}[style=pythonStyle]
a = int(input("Escribe el primer numero: "))
b = int(input("Escribe el segundo numero: "))
c = int(input("Escribe el tercer numero: "))

if a <= b and a <= c:
    menor = a
    if b <= c:
        medio, mayor = b, c
    else:
        medio, mayor = c, b
elif b <= a and b <= c:
    menor = b
    if a <= c:
        medio, mayor = a, c
    else:
        medio, mayor = c, a
else:
    menor = c
    if a <= b:
        medio, mayor = a, b
    else:
        medio, mayor = b, a

print("Orden ascendente:", menor, medio, mayor)
\end{lstlisting}

\begin{cajaSolucion}
\textbf{Alternativa con función \texttt{sorted}:}
\begin{lstlisting}[style=pythonStyle]
a = int(input("Escribe el primer numero: "))
b = int(input("Escribe el segundo numero: "))
c = int(input("Escribe el tercer numero: "))

ordenados = sorted([a, b, c])
print("Orden ascendente:", ordenados[0], ordenados[1], ordenados[2])
\end{lstlisting}
\end{cajaSolucion}
\end{cajaSolucion}

\begin{cajaSolucion}
\textbf{Reto 5}
\begin{lstlisting}[style=pythonStyle]
edad = int(input("Escribe tu edad: "))

if edad < 0:
    print("Edad invalida")
elif edad <= 2:
    print("Bebe")
elif edad <= 11:
    print("Nino(a)")
elif edad <= 17:
    print("Adolescente")
elif edad <= 59:
    print("Adulto")
else:
    print("Adulto mayor")
\end{lstlisting}

\begin{cajaNota}
Como las condiciones se evalúan en orden, \texttt{elif edad <= 2} ya implica que la edad es $\geq 0$ y $\leq 2$. No hace falta escribir \texttt{edad >= 3 and edad <= 11}, aunque también sería correcto hacerlo explícito.
\end{cajaNota}
\end{cajaSolucion}

% ============================================================
% MÓDULO 6
% ============================================================
\section{Soluciones del Módulo 6}

\begin{cajaSolucion}
\textbf{Ejercicio 21}
\begin{lstlisting}[style=pythonStyle]
for numero in range(2, 21, 2):
    print(numero)
\end{lstlisting}
\end{cajaSolucion}

\begin{cajaSolucion}
\textbf{Ejercicio 22}
\begin{lstlisting}[style=pythonStyle]
n = int(input("Escribe un numero: "))

for i in range(1, 11):
    print(n, "x", i, "=", n * i)
\end{lstlisting}
\end{cajaSolucion}

\begin{cajaSolucion}
\textbf{Ejercicio 23}
\begin{lstlisting}[style=pythonStyle]
suma = 0

for i in range(5):
    nota = float(input("Escribe la calificacion " + str(i + 1) + ": "))
    suma = suma + nota

promedio = suma / 5
print("El promedio es:", round(promedio, 2))
\end{lstlisting}
\end{cajaSolucion}

\begin{cajaSolucion}
\textbf{Ejercicio 24}
\begin{lstlisting}[style=pythonStyle]
for numero in range(10, 0, -1):
    print(numero)

print("Despegue!")
\end{lstlisting}

\begin{cajaNota}
En \texttt{range(10, 0, -1)} el paso negativo hace la cuenta regresiva: 10, 9, ..., 1. El 0 no se incluye porque el final es exclusivo.
\end{cajaNota}
\end{cajaSolucion}

\begin{cajaSolucion}
\textbf{Ejercicio 25}
\begin{lstlisting}[style=pythonStyle]
suma = 0
numero = int(input("Escribe un numero (0 para terminar): "))

while numero != 0:
    suma = suma + numero
    numero = int(input("Escribe otro numero (0 para terminar): "))

print("La suma de los numeros es:", suma)
\end{lstlisting}
\end{cajaSolucion}

\begin{cajaSolucion}
\textbf{Reto 6}
\begin{lstlisting}[style=pythonStyle]
n = int(input("Escribe un numero para calcular su factorial: "))

factorial = 1
for i in range(1, n + 1):
    factorial = factorial * i

print("El factorial de", n, "es:", factorial)
\end{lstlisting}

\textbf{Verificaciones:} $5! = 120$ y $7! = 5040$.
\end{cajaSolucion}

% ============================================================
% MÓDULO 7
% ============================================================
\section{Soluciones del Módulo 7}

\begin{cajaSolucion}
\textbf{Ejercicio 26}
\begin{lstlisting}[style=pythonStyle]
palabra = input("Escribe una palabra: ")

print("En mayusculas:", palabra.upper())
print("En minusculas:", palabra.lower())
print("Cantidad de letras:", len(palabra))
print("Primera letra:", palabra[0])
print("Ultima letra:", palabra[-1])
\end{lstlisting}
\end{cajaSolucion}

\begin{cajaSolucion}
\textbf{Ejercicio 27}
\begin{lstlisting}[style=pythonStyle]
nombre = input("Escribe tu nombre completo: ")

print(nombre.title())
\end{lstlisting}
\end{cajaSolucion}

\begin{cajaSolucion}
\textbf{Ejercicio 28}
\begin{lstlisting}[style=pythonStyle]
frase = input("Escribe una frase: ")

# Convertimos a minusculas para que A y a cuenten igual
frase_minuscula = frase.lower()
cantidad = frase_minuscula.count("a")

print("La letra 'a' aparece", cantidad, "veces")
\end{lstlisting}
\end{cajaSolucion}

\begin{cajaSolucion}
\textbf{Ejercicio 29}
\begin{lstlisting}[style=pythonStyle]
palabra = input("Escribe una palabra: ")

invertida = palabra[::-1]

if palabra == invertida:
    print(palabra, "es un palindromo")
else:
    print(palabra, "no es un palindromo")
\end{lstlisting}
\end{cajaSolucion}

\begin{cajaSolucion}
\textbf{Ejercicio 30}
\begin{lstlisting}[style=pythonStyle]
frase = input("Escribe una frase: ")

palabras = frase.split(" ")
palabras.reverse()

resultado = " ".join(palabras)
print("Frase invertida:", resultado)
\end{lstlisting}
\end{cajaSolucion}

\begin{cajaSolucion}
\textbf{Reto 7}
\begin{lstlisting}[style=pythonStyle]
frase = input("Escribe un mensaje para cifrar: ")

cifrado = ""

for letra in frase.lower():
    if letra.isalpha():        # solo ciframos letras
        codigo = ord(letra)    # codigo ASCII de la letra
        codigo = codigo + 3    # desplazamiento de 3
        # Si pasamos de 'z' (122), volvemos al inicio del alfabeto
        if codigo > ord("z"):
            codigo = codigo - 26
        cifrado = cifrado + chr(codigo)
    else:
        cifrado = cifrado + letra   # espacios y signos no cambian

print("Mensaje cifrado:", cifrado)
\end{lstlisting}

\begin{cajaNota}
Este es el clásico \textbf{cifrado César}. Con desplazamiento 3: ``hola'' se convierte en ``krod''. Para descifrar, basta restar 3 en lugar de sumar.
\end{cajaNota}
\end{cajaSolucion}

% ============================================================
% MÓDULO 8
% ============================================================
\section{Soluciones del Módulo 8}

\begin{cajaSolucion}
\textbf{Ejercicio 31}
\begin{lstlisting}[style=pythonStyle]
colores = ["rojo", "verde", "azul", "amarillo", "negro"]

print("Primer color:", colores[0])
print("Ultimo color:", colores[-1])

colores.append("blanco")
colores[1] = "naranja"

print("Lista final:", colores)
print("Tamano de la lista:", len(colores))
\end{lstlisting}
\end{cajaSolucion}

\begin{cajaSolucion}
\textbf{Ejercicio 32}
\begin{lstlisting}[style=pythonStyle]
numeros = []

for i in range(5):
    n = float(input("Escribe el numero " + str(i + 1) + ": "))
    numeros.append(n)

suma = sum(numeros)
promedio = suma / len(numeros)

print("La suma es:", suma)
print("El promedio es:", round(promedio, 2))
\end{lstlisting}
\end{cajaSolucion}

\begin{cajaSolucion}
\textbf{Ejercicio 33}
\begin{lstlisting}[style=pythonStyle]
notas = []

for i in range(3):
    n = float(input("Escribe la nota " + str(i + 1) + " (0 a 20): "))
    notas.append(n)

promedio = sum(notas) / len(notas)

print("Promedio:", round(promedio, 2))
print("Nota mas alta:", max(notas))
print("Nota mas baja:", min(notas))

notas.reverse()
print("Notas de mayor a menor:", notas)
\end{lstlisting}

\begin{cajaNota}
Alternativa para ordenar de mayor a menor: \texttt{sorted(notas, reverse=True)}.
\end{cajaNota}
\end{cajaSolucion}

\begin{cajaSolucion}
\textbf{Ejercicio 34}
\begin{lstlisting}[style=pythonStyle]
numeros = [4, 7, 2, 9, 10, 5]
pares = 0
impares = 0

for n in numeros:
    if n % 2 == 0:
        pares = pares + 1
    else:
        impares = impares + 1

print(pares, "pares y", impares, "impares")
\end{lstlisting}
\end{cajaSolucion}

\begin{cajaSolucion}
\textbf{Ejercicio 35}
\begin{lstlisting}[style=pythonStyle]
nombres = []
nombre = input("Escribe un nombre (o 'fin' para terminar): ")

while nombre != "fin":
    nombres.append(nombre)
    nombre = input("Escribe otro nombre (o 'fin' para terminar): ")

print("Se ingresaron", len(nombres), "nombres")

# Nombre mas largo
nombre_largo = nombres[0]
for n in nombres:
    if len(n) > len(nombre_largo):
        nombre_largo = n
print("El nombre mas largo es:", nombre_largo)

# Lista ordenada alfabeticamente
nombres.sort()
print("Ordenados:", nombres)
\end{lstlisting}
\end{cajaSolucion}

\begin{cajaSolucion}
\textbf{Reto 8}
\begin{lstlisting}[style=pythonStyle]
limite = int(input("Hasta que numero quieres generar la serie? "))

fibonacci = [0, 1]

while fibonacci[-1] + fibonacci[-2] <= limite:
    sgte = fibonacci[-1] + fibonacci[-2]
    fibonacci.append(sgte)

print("Serie de Fibonacci:", fibonacci)
\end{lstlisting}

\textbf{Ejemplo con límite 50:} \texttt{[0, 1, 1, 2, 3, 5, 8, 13, 21, 34]}.
\end{cajaSolucion}

% ============================================================
% MÓDULO 9
% ============================================================
\section{Soluciones del Módulo 9}

\begin{cajaSolucion}
\textbf{Ejercicio 36}
\begin{lstlisting}[style=pythonStyle]
libro = {
    "titulo": "Cien anios de soledad",
    "autor": "Gabriel Garcia Marquez",
    "anio": 1967
}

print("Titulo:", libro["titulo"])

libro["anio"] = 1967    # el anio ya es correcto
print("Anio:", libro["anio"])

libro["genero"] = "Novela"
print("Claves:", list(libro.keys()))
print("Valores:", list(libro.values()))
\end{lstlisting}
\end{cajaSolucion}

\begin{cajaSolucion}
\textbf{Ejercicio 37}
\begin{lstlisting}[style=pythonStyle]
estaciones = ("primavera", "verano", "otonio", "invierno")

print("Primera estacion:", estaciones[0])
print("Ultima estacion:", estaciones[-1])
print("Cantidad de estaciones:", len(estaciones))

# Las tuplas son inmutables: no se puede modificar un elemento
# estaciones[0] = "otono"   # esto generaria un error
\end{lstlisting}
\end{cajaSolucion}

\begin{cajaSolucion}
\textbf{Ejercicio 38}
\begin{lstlisting}[style=pythonStyle]
capitales = {}

for i in range(3):
    pais = input("Escribe el nombre de un pais: ")
    capital = input("Escribe la capital de " + pais + ": ")
    capitales[pais] = capital

for pais, capital in capitales.items():
    print("La capital de", pais, "es", capital)
\end{lstlisting}
\end{cajaSolucion}

\begin{cajaSolucion}
\textbf{Ejercicio 39}
\begin{lstlisting}[style=pythonStyle]
notas = {"Ana": 18, "Luis": 14, "Marta": 20, "Pedro": 11}

promedio = sum(notas.values()) / len(notas)
print("Promedio:", round(promedio, 2))

nota_max = max(notas.values())
nota_min = min(notas.values())

for estudiante, nota in notas.items():
    if nota == nota_max:
        print("La nota mas alta es de", estudiante, "con", nota)
    if nota == nota_min:
        print("La nota mas baja es de", estudiante, "con", nota)
\end{lstlisting}
\end{cajaSolucion}

\begin{cajaSolucion}
\textbf{Ejercicio 40}
\begin{lstlisting}[style=pythonStyle]
palabra = input("Escribe una palabra: ")

conteo = {}

for letra in palabra:
    if letra in conteo:
        conteo[letra] = conteo[letra] + 1
    else:
        conteo[letra] = 1

print(conteo)
\end{lstlisting}

\textbf{Con ``banana'':} \texttt{\{'b': 1, 'a': 3, 'n': 2\}}.

\begin{cajaNota}
Alternativa con el método \texttt{get}: \texttt{conteo[letra] = conteo.get(letra, 0) + 1}. Hace lo mismo en una sola línea.
\end{cajaNota}
\end{cajaSolucion}

\begin{cajaSolucion}
\textbf{Reto 9}
\begin{lstlisting}[style=pythonStyle]
agenda = {}

while True:
    print("\n--- AGENDA TELEFONICA ---")
    print("1. Agregar contacto")
    print("2. Buscar contacto")
    print("3. Eliminar contacto")
    print("4. Mostrar todos")
    print("5. Salir")

    opcion = input("Elige una opcion: ")

    if opcion == "1":
        nombre = input("Nombre del contacto: ")
        numero = input("Numero de telefono: ")
        agenda[nombre] = numero
        print("Contacto agregado correctamente")

    elif opcion == "2":
        nombre = input("Que contacto buscas? ")
        if nombre in agenda:
            print("El numero de", nombre, "es", agenda[nombre])
        else:
            print("Ese contacto no existe")

    elif opcion == "3":
        nombre = input("Que contacto quieres eliminar? ")
        if nombre in agenda:
            del agenda[nombre]
            print("Contacto eliminado")
        else:
            print("Ese contacto no existe")

    elif opcion == "4":
        if len(agenda) == 0:
            print("La agenda esta vacia")
        else:
            for nombre, numero in agenda.items():
                print(nombre, "->", numero)

    elif opcion == "5":
        print("Hasta luego!")
        break

    else:
        print("Opcion no valida, intenta de nuevo")
\end{lstlisting}
\end{cajaSolucion}

% ============================================================
% MÓDULO 10
% ============================================================
\section{Soluciones del Módulo 10}

\begin{cajaSolucion}
\textbf{Ejercicio 41}
\begin{lstlisting}[style=pythonStyle]
def dibujar_linea():
    print("*" * 20)

dibujar_linea()
dibujar_linea()
dibujar_linea()
\end{lstlisting}
\end{cajaSolucion}

\begin{cajaSolucion}
\textbf{Ejercicio 42}
\begin{lstlisting}[style=pythonStyle]
def cuadrado(n):
    return n ** 2

numero = int(input("Escribe un numero: "))
print("El cuadrado de", numero, "es:", cuadrado(numero))
\end{lstlisting}
\end{cajaSolucion}

\begin{cajaSolucion}
\textbf{Ejercicio 43}
\begin{lstlisting}[style=pythonStyle]
def es_par(n):
    return n % 2 == 0

contador_pares = 0

for i in range(5):
    numero = int(input("Escribe el numero " + str(i + 1) + ": "))
    if es_par(numero):
        contador_pares = contador_pares + 1

print("Hay", contador_pares, "numeros pares")
\end{lstlisting}
\end{cajaSolucion}

\begin{cajaSolucion}
\textbf{Ejercicio 44}
\begin{lstlisting}[style=pythonStyle]
def promedio(notas):
    return sum(notas) / len(notas)

mis_notas = []
for i in range(4):
    n = float(input("Escribe la nota " + str(i + 1) + ": "))
    mis_notas.append(n)

print("El promedio es:", round(promedio(mis_notas), 2))
\end{lstlisting}
\end{cajaSolucion}

\begin{cajaSolucion}
\textbf{Ejercicio 45}
\begin{lstlisting}[style=pythonStyle]
def invertir_texto(texto):
    return texto[::-1]

frase = input("Escribe una frase: ")
print("Frase invertida:", invertir_texto(frase))
\end{lstlisting}
\end{cajaSolucion}

\begin{cajaSolucion}
\textbf{Reto 10}
\begin{lstlisting}[style=pythonStyle]
def sumar(a, b):
    return a + b

def restar(a, b):
    return a - b

def multiplicar(a, b):
    return a * b

def dividir(a, b):
    if b == 0:
        return "No se puede dividir entre cero"
    return a / b

while True:
    print("\n--- CALCULADORA ---")
    print("1. Sumar")
    print("2. Restar")
    print("3. Multiplicar")
    print("4. Dividir")
    print("5. Salir")

    opcion = input("Elige una operacion: ")

    if opcion == "5":
        print("Hasta luego!")
        break

    a = float(input("Primer numero: "))
    b = float(input("Segundo numero: "))

    if opcion == "1":
        print("Resultado:", sumar(a, b))
    elif opcion == "2":
        print("Resultado:", restar(a, b))
    elif opcion == "3":
        print("Resultado:", multiplicar(a, b))
    elif opcion == "4":
        print("Resultado:", dividir(a, b))
    else:
        print("Opcion no valida")
\end{lstlisting}
\end{cajaSolucion}

% ============================================================
% MÓDULO 11
% ============================================================
\section{Soluciones del Módulo 11}

\begin{cajaSolucion}
\textbf{Ejercicio 46}
\begin{lstlisting}[style=pythonStyle]
import math

print("Raiz cuadrada de 289:", math.sqrt(289))
print("Pi redondeado:", round(math.pi, 4))
print("Factorial de 10:", math.factorial(10))
\end{lstlisting}
\end{cajaSolucion}

\begin{cajaSolucion}
\textbf{Ejercicio 47}
\begin{lstlisting}[style=pythonStyle]
import random

for i in range(3):
    input("Presiona Enter para lanzar el dado...")
    resultado = random.randint(1, 6)
    print("Lanzamiento", i + 1, ":", resultado)
\end{lstlisting}
\end{cajaSolucion}

\begin{cajaSolucion}
\textbf{Ejercicio 48}
\begin{lstlisting}[style=pythonStyle]
import random

secreto = random.randint(1, 50)
intentos = 0
adivinado = False

while not adivinado:
    numero = int(input("Adivina el numero (1 a 50): "))
    intentos = intentos + 1

    if numero < secreto:
        print("El numero es mayor")
    elif numero > secreto:
        print("El numero es menor")
    else:
        print("Correcto! Lo lograste en", intentos, "intentos")
        adivinado = True
\end{lstlisting}
\end{cajaSolucion}

\begin{cajaSolucion}
\textbf{Ejercicio 49}
\begin{lstlisting}[style=pythonStyle]
import math

a = float(input("Escribe el primer cateto: "))
b = float(input("Escribe el segundo cateto: "))

hipotenusa = math.sqrt(math.pow(a, 2) + math.pow(b, 2))
print("La hipotenusa es:", round(hipotenusa, 2))
\end{lstlisting}
\end{cajaSolucion}

\begin{cajaSolucion}
\textbf{Ejercicio 50}
\begin{lstlisting}[style=pythonStyle]
import random

respuestas = ["Si", "No", "Quizas", "Sin duda", "Claro que no"]

pregunta = input("Hazme una pregunta: ")
print("Respuesta:", random.choice(respuestas))
\end{lstlisting}
\end{cajaSolucion}

\begin{cajaSolucion}
\textbf{Reto 11}
\begin{lstlisting}[style=pythonStyle]
import random

marcador_usuario = 0
marcador_maquina = 0
opciones = ["piedra", "papel", "tijera"]

for ronda in range(1, 6):
    print("\n--- Ronda", ronda, "---")
    usuario = input("Elige piedra, papel o tijera: ").lower()
    maquina = random.choice(opciones)
    print("La maquina eligio:", maquina)

    if usuario == maquina:
        print("Empate en esta ronda")
    elif (usuario == "piedra" and maquina == "tijera") or \
         (usuario == "papel" and maquina == "piedra") or \
         (usuario == "tijera" and maquina == "papel"):
        print("Ganaste esta ronda!")
        marcador_usuario = marcador_usuario + 1
    else:
        print("La maquina gano esta ronda")
        marcador_maquina = marcador_maquina + 1

print("\n--- RESULTADO FINAL ---")
print("Usuario:", marcador_usuario, "- Maquina:", marcador_maquina)

if marcador_usuario > marcador_maquina:
    print("Ganaste el juego!")
elif marcador_maquina > marcador_usuario:
    print("La maquina gano el juego")
else:
    print("Empate general")
\end{lstlisting}

\begin{cajaNota}
En Python la barra invertida al final de una línea (\texttt{\textbackslash}) permite continuar una expresión larga en la siguiente línea, una técnica muy útil para condiciones con muchas partes.
\end{cajaNota}
\end{cajaSolucion}

% ============================================================
% PROYECTO FINAL
% ============================================================
\section{Sugerencias para el Proyecto Final}

Los proyectos finales son libres; a continuación se sugieren soluciones de referencia para el docente.

\begin{cajaSolucion}
\textbf{Proyecto E (solución de referencia): Registro de notas}

\begin{lstlisting}[style=pythonStyle]
# Registro de notas: cada estudiante tiene una lista de notas
registro = {}

def registrar_estudiante(nombre):
    if nombre not in registro:
        registro[nombre] = []

def agregar_nota(nombre, nota):
    if nombre in registro:
        registro[nombre].append(nota)

def promedio_lista(lista):
    return sum(lista) / len(lista)

def reporte():
    for estudiante, notas in registro.items():
        promedio = promedio_lista(notas)
        print(estudiante, "-> notas:", notas,
              "-> promedio:", round(promedio, 2))

while True:
    print("\n--- REGISTRO DE NOTAS ---")
    print("1. Registrar estudiante")
    print("2. Agregar nota")
    print("3. Notas de un estudiante")
    print("4. Reporte general y promedio del curso")
    print("5. Salir")

    opcion = input("Elige una opcion: ")

    if opcion == "1":
        nombre = input("Nombre del estudiante: ")
        registrar_estudiante(nombre)
        print("Estudiante registrado")

    elif opcion == "2":
        nombre = input("Nombre del estudiante: ")
        if nombre not in registro:
            print("Ese estudiante no esta registrado")
        else:
            nota = float(input("Escribe la nota: "))
            agregar_nota(nombre, nota)
            print("Nota agregada")

    elif opcion == "3":
        nombre = input("Nombre del estudiante: ")
        if nombre in registro:
            print(nombre, "->", registro[nombre])
        else:
            print("Ese estudiante no existe")

    elif opcion == "4":
        if len(registro) == 0:
            print("No hay estudiantes registrados")
        else:
            reporte()
            todas = []
            for notas in registro.values():
                todas = todas + notas
            print("Promedio del curso:",
                  round(sum(todas) / len(todas), 2))

    elif opcion == "5":
        print("Hasta luego!")
        break

    else:
        print("Opcion no valida")
\end{lstlisting}
\end{cajaSolucion}

\end{document}